% ===== phụ lục (không giới hạn trang) =====
\textbf{T1 --- khử bóng} (\texttt{t1\_background.py}; mẫu chỉ có bóng: 10 valid, 15 test; PSNR / SSIM so với ảnh sạch xám).
\begin{table}[H]
\centering\small
\begin{tabular}{lrrr}
\toprule
Ước lượng nền & SE / cửa sổ (chọn trên valid) & PSNR test (dB) & SSIM test \\
\midrule
Không khử & --- & 11,67 & 0,886 \\
Trung vị cỡ lớn (C3 S73) & 101 & 35,24 & 0,995 \\
\textbf{Closing ảnh xám} (C3 S74, C7 S35) & 81 & \textbf{47,91} & \textbf{0,9993} \\
Closing $+$ trung vị 21 & 81 & 48,10 & 0,9994 \\
\bottomrule
\end{tabular}
\end{table}
Closing $+$ trung vị nhỉnh $0{,}2$~dB nhưng chậm hơn; giữ Closing 81 (trên valid Closing đơn thuần cao hơn: 52,79 so với 52,64~dB).

\textbf{T4 --- dấu đỏ} (\texttt{t4\_stamp.py}; 40 mẫu có dấu, val $+$ test).
\begin{table}[H]
\centering\small
\begin{tabular}{lrrrr}
\toprule
Mức dấu & IoU mask đỏ HSV & IoU nét dấu tìm được & Phủ nét dấu & Tìm được ($\ge80\%$) \\
\midrule
1 & 0,534 & 0,534 & 0,950 & 100\% \\
2 & 0,741 & 0,679 & 0,896 & 92\% \\
3 & 0,809 & 0,746 & 0,903 & 92\% \\
\textbf{Tất cả} & 0,690 & 0,650 & 0,917 & 95\% \\
\bottomrule
\end{tabular}
\end{table}
PSNR trong vùng dấu sau khi xóa: tô màu giấy \SI{13.9}{dB}, Telea \SI{13.5}{dB}, Navier--Stokes \SI{13.4}{dB}, không xóa \SI{13.0}{dB}
--- ba cách xóa gần nhau. Ảnh hưởng lên OCR đo ở T8: xóa dấu / bút giảm CER từ 0,539 xuống 0,433 (dấu) và từ 0,384 xuống 0,238 (bút).

\begin{table}[H]
\centering
\caption{T2-CER: CER toàn tập khi OCR từng dòng (Tesseract 5) cắt từ ảnh thật đã nhị phân cả trang (sau khử bóng). 600 dòng VAIPE
mỗi tập; 600 cụm từ của 32 phiếu \bo{2} (phóng $\times3$). Tiêu chí chốt trước khi chạy: chọn theo VAIPE valid.}\label{tab:T2CER}
\small
\begin{tabular}{lrrr}
\toprule
Đầu vào OCR & VAIPE valid (chọn) & VAIPE test & \Bo{2} \\
\midrule
Ảnh xám đã khử bóng (không nhị phân) & 0,063 & 0,046 & 0,300 \\
\midrule
Niblack & \textbf{0,146} & \textbf{0,142} & 0,568 \\
\textbf{Sauvola} & 0,160 & 0,159 & \textbf{0,538} \\
Tam giác / Đối xứng nền & 0,256 & 0,205 / 0,217 & 0,955 / 0,991 \\
ISODATA & 0,281 & 0,259 & 0,558 \\
Otsu & 0,282 & 0,260 & 0,551 \\
Wolf & 0,296 & 0,271 & 0,627 \\
\bottomrule
\end{tabular}
\\[0.3em]{\footnotesize Script: \texttt{t2cer\_binarize\_ocr.py}}
\end{table}

\textbf{T5 --- giữ chữ đỏ trên phiếu xét nghiệm thật} (\texttt{t5\_red\_text.py}; định tính trên 5 phiếu \bo{4} chọn bằng mắt
vì có mực đỏ thật: 003, 004, 028, 041, 048). Mỗi phiếu chạy (1) $\rightarrow$ (2) nhánh phiếu xét nghiệm; ảnh trước / sau ở
\texttt{results/t5/} (có thông tin bệnh nhân, không đưa vào repo).
\begin{table}[H]
\centering\small
\begin{tabularx}{\linewidth}{lrrX}
\toprule
Phiếu & Pixel đỏ & Bị coi là dấu & Nhận xét bằng mắt \\
\midrule
003 & \num{1349} & 0 & Chữ đỏ cảnh báo ``10.29'', ``Dương tính'' giữ nguyên và còn rõ trong ảnh nhị phân (Hình~\ref{fig:T5}) \\
004 & \num{16977} & 0 & Giấy hồng: đường kẻ bảng bị mask đỏ bắt nhầm (dương tính giả), không có gì bị xóa \\
028 & \num{90131} & 0 & Khối chữ đầu phiếu (người lấy / nhận mẫu) bị mask đỏ đánh dấu là chữ đỏ và được giữ; mép tay / giấy hồng ở góc là dương tính giả \\
041 & \num{26052} & 0 & Logo đỏ in chìm sau bảng: chỉ 2 mảng đậm nhất vượt ngưỡng $S$, Sauvola loại phần còn lại khỏi ảnh nhị phân \\
048 & \num{7828} & 0 & Dấu đỏ \emph{chữ nhật} của phòng xét nghiệm: độ tròn thấp nên được giữ như chữ đỏ, đọc được trong ảnh nhị phân \\
\bottomrule
\end{tabularx}
\end{table}
Không phiếu nào có thành phần đỏ bị R4 coi là dấu, nên kể cả khi (nhầm) chạy nhánh xóa dấu của đơn thuốc thì cũng không mất chữ
đỏ nào; số phiếu có chữ đỏ bị xóa nhầm: \textbf{0/5}. Hai ca lỗi không thuộc bước tách mực: phiếu 003 bị bước (1) chọn khung bảng
làm ``tờ giấy'' nên mất phần đầu phiếu; phiếu 048 có vùng bóng tay đè lên cột kết quả dòng 2--7, Sauvola làm mất các số in mờ
trong vùng này. Mask đỏ HSV có dương tính giả trên giấy hồng / vàng và da tay, vô hại với phiếu xét nghiệm (không xóa gì) nhưng
là rủi ro nếu ảnh đơn thuốc chụp trên nền như vậy.
\begin{figure}[H]\centering\includegraphics[width=\linewidth]{cleaning/T5.jpg}
\caption{T5, phiếu 003 (vùng kết quả, không có thông tin cá nhân): ảnh đã nắn $|$ chữ đỏ được giữ (xanh lá) $|$ ảnh nhị phân sau (2).}\label{fig:T5}\end{figure}

\textbf{Bảng tham số} (KeHoach.md mục 6.3).
\begin{table}[H]
\centering\small
\begin{tabularx}{\linewidth}{lllX}
\toprule
Tham số & Khoảng thử & Chọn & Lý do \\
\midrule
SE ước lượng nền & 31, 51, 81, 101 & 81 & PSNR lớn nhất trên valid (T1) \\
Phương pháp nhị phân & 7 cách & Sauvola (cho tách vùng) & Bền nhất khi không khử bóng (T2, T3). Theo T2-CER, Niblack có CER dòng thấp nhất nhưng OCR trên ảnh xám còn tốt hơn, nên nhị phân chỉ dùng tách vùng \\
Cửa sổ Sauvola ($\times h$) & 1,5--4 & 2,5 & F-measure (bỏ dải viền) lớn nhất trên valid \\
$k$ Sauvola & 0,1--0,5 & 0,3 & như trên \\
Opening hậu xử lý & tắt, $2\times2$ & tắt & Opening xóa nét chữ 1--2~px (C7 S14) \\
Ngưỡng diện tích nhiễu & $(0{,}4h)^2$, $\ge3$~px & theo $h$ & Ngưỡng cố định 10~px xóa sạch chữ ảnh 640~px \\
Độ tròn dấu; diện tích & 0,6; $\ge0{,}2\%$ trang & --- & Phát hiện 95\% dấu (T4) \\
\bottomrule
\end{tabularx}
\end{table}
